\section{Preguntas de análisis}
\label{sec:questions}

Las preguntas (Tabla~\ref{tab:questions}) van de la descripción del problema a su relación con la salud, el
agua y el clima. Todas, salvo la primera, integran al menos dos datasets.

\begin{table}[htbp]
\centering
\caption{Preguntas de análisis y datasets que usa cada una.}
\label{tab:questions}
\small
\begin{tabularx}{\textwidth}{@{}l X >{\raggedright\arraybackslash}p{3.8cm}@{}}
\toprule
 & Pregunta & Datasets \\
\midrule
Q1 & ¿Cómo evolucionó el dengue entre 2000 y 2025 y qué años fueron epidémicos? & dengue \\
Q2 & ¿Qué distritos tuvieron la mayor incidencia (casos por 100,000 habitantes) en 2022–2025? & dengue, población \\
Q3 & ¿Qué grupos de edad y qué sexo tienen la mayor incidencia? & dengue, población \\
Q4 & ¿La proporción de casos con signos de alarma o graves se relaciona con la cobertura de establecimientos de
     salud? & dengue, establecimientos, población \\
Q5 & ¿Los distritos en emergencia hídrica tuvieron mayor incidencia? ¿Influye la continuidad del servicio de
     agua? & dengue, población, emergencias \\
Q6 & ¿Cuántas semanas después de la lluvia aumentan los casos, y es distinto en la costa y en la selva? &
     dengue, estaciones \\
\bottomrule
\end{tabularx}
\end{table}
